\documentclass[exercice]{thedocument}%
\usepackage{texmaths-tikz}%
\startdatakeys
\source{}
\cycle{}
\competencesDuSocle{}
\connaissancesRequises{}
\competencesTravaillees{}
\enddatakeys
\begin{document}%
\begin{center}\footnotesize
    \drawTikzOnPic[style={/utils/exec={\normalfont\bfseries}}]{%
        \pdf[scale=1.15]{pythagore_architecture}
    }{%
        \node[anchor=center] at (0.21,0.12){\<L1;};
        \node[anchor=center] at (0.78,0.12){\<L2;};
        \node[anchor=center,rotate=-90] at (0.41,0.3){\<L3;};
        \node[anchor=center,rotate=90] at (0.57,0.3){\<L3;};
        \node[anchor=center] at (0.49,0.7){\<L4;};
    }
\end{center}
\begin{enumerate}
    \item Déterminer une valeur approchée, au centimètre près,
    de la longueur du tablier de ce pont.
    \item Ce pont a une largeur de \<largeurPont;.\par
    En déduire un valeur approchée de l'aire du plancher constituant
    le sol du pont.
\end{enumerate}
\end{document}
